\chapter{L’imputation des données manquantes : un problème transversal et un champ de recherche dynamique}

\section{Les données manquantes : un problème connu et commun à travers les domaines}

% =====================================================================
% 1. DÉFINITION ET CAUSES
% =====================================================================

On parle de donnée manquante lorsqu’une valeur n’a pas été enregistrée pour une variable ou pour un individu \cite{Zahmatkesh_2026}.
% Source : Zahmatkesh, §2.2 « Missing Data », p. 6, 1re phrase.
% Citation complète dans ta mind map (branche « What are missing values? »).
Ce phénomène est impossible à éviter complètement. 
Même lorsque le système de collecte est bien conçu et que les capteurs sont de bonne qualité, des valeurs finissent toujours par manquer \cite{Adhikari_2023}. 
Les causes sont très variées, 
% Source 1 : Adhikari, §1 « Introduction », p. 2, vers la fin du 1er paragraphe :
%   « No matter how strictly the various intelligent systems are designed... »
% Source 2 : Zahmatkesh, §2.2, p. 6, 2e phrase : liste « human errors, technical failures,
%   data entry issues, data processing errors, data integration problems, or intentional omissions ».
Adhikari et al. proposent de les classer selon l’architecture d’un système de type Internet des objets (IoT), qui comporte trois couches \cite{Adhikari_2023} :
\begin{itemize}
    \vspace{2mm}
    \item \textbf{La couche physique (capteurs)} : environnement difficile, phénomènes météorologiques extrêmes, maintenance, vieillissement ou panne des capteurs, batterie ou alimentation insuffisante, attaques sur les capteurs ;
    \vspace{2mm}
    \item \textbf{La couche réseau (transmission)} : perte de paquets, problèmes de chiffrement ou de déchiffrement, équipements défaillants, attaques sur le réseau, problèmes d’alimentation ;
    \vspace{2mm}
    \item \textbf{La couche application (stockage et traitement)} : erreur humaine (suppression, oubli), non-réponse à une enquête, erreur lors de l’enregistrement ou du transfert, ou cause inconnue.
    \vspace{5mm}
\end{itemize}
% Source : Adhikari, §2.2 « IoT Architecture and Missing Data », pp. 5-6 + Fig. 1 (p. 6).
%   Le texte décrit les 3 couches (« sensing/object layer », « network layer », « application layer »).
%   Le détail des listes (perte de paquets, vieillissement, batterie...) vient surtout de la Fig. 1 et de §1/§2.1.
% \textit{DOUTE : j’ai complété certains points (« perte de paquets », « problème d’alimentation »)
%   à partir de la transcription de la Fig. 1 dans ta mind map. Vérifie sur la figure d’origine.}
% \textit{GRAPHE SUGGÉRÉ : un schéma à trois couches (capteurs, réseau, application) avec les causes
%   dans chaque couche. À redessiner toi-même plutôt que copier la Fig. 1 d’Adhikari (droits d’auteur).}


% =====================================================================
% 2. DOMAINES TOUCHÉS
% =====================================================================

On retrouve les mêmes causes dans d’autres applications; en realité, le problème touche presque tous les domaines. 
Voici quelques exemples concrets, qui concernent directement les séries temporelles multivariées:
\begin{itemize}
    \vspace{2mm}
    \item En santé, le jeu de données PhysioNet 2012, composé de séjours en soins intensifs, est très souvent utilisé comme référence. 
    Son taux global de valeurs manquantes dépasse 80~\% car elles ne sont réalisées que lors de l'intervention d'un spécialiste \cite{Luo_2018,Richter_2024,Wang_2025}.
    \vspace{2mm}
    \item Pour la qualité de l'air, le jeu de données italien PT08.S1 présente environ 34~\% de valeurs manquantes \cite{Richter_2024}.
    \vspace{2mm}
    \item Les mesures de température, de précipitations ou d’humidité proviennent surtout de réseaux de stations, souvent incomplets \cite{Alejo-Sanchez_2025}.
    \vspace{2mm}
    \item Fang et Wang montrent que des erreurs de capteurs de niveau d'essence créent des trous dans les données, et que les imputer améliore nettement la prédiction de consommation \cite{Fang_2020}.
\end{itemize} \vspace{5mm}
% Source 1 (santé) : Luo, §1, p. 2 + Fig. 2 : « The global missing rate is 80.67% » ; Richter, §I, p. 1 (148167) :
%   « Physionet 2012 data set on medical data has a missing rate of 80 percent » ; Wang, §1, p. 1 : « missing rates can exceed 80% ».
% Source 2 (air) : Richter, §I, p. 1 (148167) : PT08.S1 « has a missing rate of 34 percent ».
% Source 3 (climat) : Alejo-Sanchez, « Study areas analysis », p. 4 + Fig. 2 : 93 % des études utilisent des réseaux de stations.
% Source 4 (carburant) : Fang, §1, p. 1, 3e paragraphe : « imputing the missing values indeed helps significantly
%   the prediction of fuel consumption ».
% Source 5 (météo spatiale) : Rüdisser et al., 2026, p. 1, uniquement via ta mind map (branche « Why is it a problem? »).
% \textit{DOUTE : Rüdisser n’est pas dans les PDF que tu m’as envoyés. Je me suis appuyé uniquement sur ta mind map.
%   Je n’ai pas pu vérifier la page ni la formulation. Clé BibTeX supposée : rudisser_2026.}
% \textit{DOUTE : pour la santé, Luo parle de 4 000 séjours et Richter de 12 000 admissions. Je n’ai donc pas donné de taille de jeu de données.}
% \textit{GRAPHE SUGGÉRÉ : si tu utilises PhysioNet dans ton projet, ta propre carte de chaleur des valeurs manquantes
%   (variables × heures) serait plus parlante, sur le modèle des Fig. 1-2 de Luo.}

Le taux de valeurs manquantes peut être très élevé. 
Plusieurs auteurs indiquent qu’il peut atteindre 90~\% dans certains jeux de données \cite{Fang_2020,Richter_2024}. \\
% Source 1 : Fang, §1, p. 1, 1er paragraphe : « the missing rate can reach 90% ».
% Source 2 : Richter, §I, p. 1 (148167) : « the missing rate can reach up to 90 % ».
% \textit{DOUTE : les deux articles citent la même source d’origine (García-Laencina et al., 2015).
%   Si tu veux la citer directement, il faudra l’ajouter à ta bibliographie.}

% =====================================================================
% 3. POURQUOI IL FAUT LES TRAITER CORRECTEMENT
% =====================================================================

Il est important de traiter correctement ces données manquantes, car ne pas le faire a des conséquences directes. 
Adhikari et al. citent plusieurs risques : des résultats biaisés, une taille d’échantillon réduite, une perte d’efficacité et de puissance statistique, des intervalles de confiance moins précis, et au final des conclusions trompeuses \cite{Adhikari_2023}.
Certains modèles de Machine Learning ont aussi besoin de données complètes pour bien fonctionner \cite{Adhikari_2023}. 
Les effets se voient en aval : classification, régression et prévision sont dégradées \cite{Fang_2020,Wang_2025}. 
Dans le cas du climat, les trous compliquent le suivi et la prévision des phénomènes \cite{Alejo-Sanchez_2025}. 
De plus, les valeurs manquantes proches dans le temps ou dans l’espace peuvent contenir de l’information utile, et doivent donc être soigneusement récupérées \cite{Zahmatkesh_2026}.
% Source 1 : Adhikari, §1.1, p. 2, 3e paragraphe : « The high risk of obtaining bias results... reduction in sample size,
%   loss of efficiency... misleading results ».
% Source 2 : Adhikari, §1.1, p. 2, 4e paragraphe : « All of these techniques require complete quality data ».
% Source 3 : Fang, Abstract, p. 1 : l’absence de valeurs nuit aux applications en aval (classification, régression,
%   intégration de données séquentielles, prévision). Wang, §1, p. 1 : même idée (« accuracy and reliability of downstream analysis »).
% Source 4 : Alejo-Sanchez, Abstract, p. 1 : « complicates the monitoring and prediction of climatic phenomena ».
% Source 5 : Zahmatkesh, §1, p. 2, 2e paragraphe : « Missing data in spatial or temporal vicinity may be informative ».
%   Attention : la phrase sur le biais dans ce paragraphe concerne l’ignorance des dépendances spatio-temporelles,
%   pas les données manquantes. Je ne l’ai donc pas utilisée ici.

Cependant, remplir les trous n’est pas forcément une tâche aisée.
Le choix de la méthode adaptée dépend de la façon dont les données manquent : leur durée, leur motif et leur mécanisme d'apparition \cite{Zahmatkesh_2026}.
C’est l’objet de la section suivante.
% Source : Zahmatkesh, §1, p. 2 : choisir une méthode d’imputation « with the knowledge of missing data mechanisms and the pattern of missingness » ;
%   §2.2, p. 6 : comprendre la distribution des trous aide à choisir l’algorithme.




\section{Comprendre la répartition des données manquantes}

\subsection{Durées}

Gap classification (Adhikari et al 2023)

\subsection{Patterns}

univariate Vs multivariate:

monotone Vs non-monotone:

file matching



\subsection{Mechanisms}


Rubin (MCAR, MAR, NMAR)

\section{Méthodes d’imputation}
(préciser qu'il existe d'autres techniques que l'imputation mais quelle sont pas du tt dans notre contexte)

\subsection{single imputation}

\subsection{multiple imputation}
\subsubsection{Model based}
\subsubsection{Imputation based}

\subsection{Méthodes statistiques conventionnelles}

\subsection{Méthode de Machine Learning (shallow learning)}

Préciser également la différence entre les modèles supervisés et non supervisés
\subsubsection{Classification}

\subsubsection{Clustering}



\subsection{Méthodes de Deep Learning}

Bien préciser que la première différence de méthodes est lié à un concept déterministe vs stochastic
\subsection{Méthodes prédictives}

\subsubsection{architectures principales}

RNN, CNN, GNN, Attention based

\subsubsection*{modèles appliqués à l'imputation}

\subsection{Méthodes génératives}

\subsubsection{architectures principales}

VAE, GAN, Diffusion
\subsubsection*{modèles appliqués à l'imputation}


\section{Évaluation des performances}
différentes métriques (Évaluation directe, correlation, mesure du biais, conservation de la distribution)

\section{PyPOTS}

en quoi pypots nous est utile